\chapter{Requirement Analysis}
\label{chap:requirement-analysis}

\section{Stakeholder Analysis}
\label{section:stakeholder-analysis}

Stakeholders are individuals or groups who have an interest in the system and may influence or be affected by its outcomes. 
For this project, the key stakeholders are identified as follows:

\begin{enumerate}[leftmargin=80pt]
    \item Deaf Users (Primary Users): The deaf and hard-of-hearing individuals who rely on the system to gain awareness of environmental sounds through visual representation.
    \item General Users (Secondary Users): The general individuals without hearing impairments who may use the application for exploration, learning, or in situations where audio perception is limited.
    \item Researchers and Educators: The researchers and educators who may use the application as a tool to study sound behavior, accessibility solutions, or human-computer interaction.
    \item Assistive Technology Stakeholders: Individuals and organizations involved in assistive technology, including accessibility experts and other stakeholders, who may evaluate or adopt the system in the future.
    \item Developers: The developers who are responsible for designing, implementing, and maintaining the system, ensuring performance, usability, and extensibility.
\end{enumerate}

\section{User Stories and Use Cases}
\label{section:user-stories-use-cases}

\subsection{User Stories}
\begin{enumerate}[leftmargin=80pt]

\item \textbf{Deaf and Hard-of-Hearing Individuals (Primary Users)}
\begin{itemize}
    \item As a deaf user, I want to see what type of sound is occurring around me so that I can understand my environment.
    \item As a deaf user, I want to know the estimated direction of a sound source so that I can locate its source.
    \item As a deaf user, I want to know the estimated distance of a sound source so that I can locate its source.
    \item As a deaf user, I want real-time updates of sound activity so that I can stay aware of changes in my surroundings.
    \item As a deaf user, I want a simple visual interface so that I can quickly interpret sound information without confusion.
\end{itemize}

\item \textbf{General Users (Secondary Users)}
\begin{itemize}
    \item As a general user, I want to explore environmental sounds visually so that I can better understand my surroundings.
    \item As a general user, I want to use the application in various difficult environments so that I can interpret sounds more clearly.
\end{itemize}

\item \textbf{Researchers and Educators}
\begin{itemize}
    \item As a researcher, I want to analyze sound data so that I can study sound behavior and interaction.
    \item As an educator, I want to demonstrate sound concepts visually so that students can better understand them.
    \item As a researcher, I want access to categorized sound data so that I can conduct experiments or evaluations.
\end{itemize}

\item \textbf{Assistive Technology Stakeholders}
\begin{itemize}
    \item As an accessibility expert, I want to evaluate the usability of the system so that I can assess its effectiveness.
    \item As a stakeholder, I want to check whether the system complies with accessibility standards so that it is usable by a wide audience.
\end{itemize}

\item \textbf{Developers}
\begin{itemize}
    \item As a developer, I want the system to be modular so that it can be easily maintained and extended.
    \item As a developer, I want the system to support efficient real-time processing so that it performs smoothly on mobile devices.
    \item As a developer, I want to integrate machine learning models so that sound classification is accurate and reliable.
    \item As a developer, I want a clear and well-defined system architecture so that future development and enhancements are easier.
\end{itemize}

\end{enumerate}

\subsection{Use Cases}

The use cases are derived from the user stories by grouping related functionalities into higher-level system interactions.

\begin{enumerate}[leftmargin=80pt]

\item \textbf{Detect and Visualize Sound (Fully-Dressed)}

Actor: User

Description: The system captures environmental audio and displays it visually.

Main Flow:
\begin{enumerate}[leftmargin=80pt]
    \item The user opens the application.
    \item The system activates the microphone.
    \item Environmental sound is captured in real-time.
    \item The system processes the sound (classification, direction, and distance).
    \item The results are displayed on the visual map interface.
\end{enumerate}

Alternative Flow:
\begin{enumerate}[leftmargin=80pt]
    \item If no sound is detected, the interface remains idle or shows a neutral state.
\end{enumerate}

\item \textbf{Identify Sound Type (Casual)}

The system analyzes incoming audio and categorizes it into predefined types such as speech, vehicle, alarm, or ambient noise. The identified type is displayed using icons or labels.

\item \textbf{View Estimated Sound Direction (Casual)}

The system analyzes incoming audio and estimates the direction of a sound source using cardinal and intermediate directions such as east, north, or east-west. The estimated direction is displayed using spatial positioning on the interface.

\item \textbf{View Estimated Sound Distance (Casual)}

The system analyzes incoming audio and estimates the distance of a sound source using distance levels such as near, far, or medium. The estimated distance is displayed using spatial positioning on the interface.

\item \textbf{Explore Sound Map (Brief)}

The user observes and interacts with the visual map to understand surrounding sounds.

\item \textbf{Analyze Sound Data (Brief)}

Researchers analyze collected sound data through visual representations to study sound behavior and interaction.

\item \textbf{Access Categorized Data (Brief)}

The system provides categorized sound data that can be used for research, evaluation, or experimentation.

\item \textbf{Evaluate System Usability (Brief)}

Accessibility experts evaluate the usability of the system to assess its effectiveness for target users.

\item \textbf{Check Accessibility Compliance (Brief)}

The system is assessed against accessibility standards to ensure it is usable by a wide audience.

\item \textbf{Maintain and Improve System (Brief)}

Developers maintain the system, optimize performance, and improve machine learning models to ensure accuracy and efficiency.

\end{enumerate}

\section{Use Case Diagrams}
\label{section:use-case-diagrams}

\subsection{Use Case Diagram}

\begin{figure}[h]
    \centering
    \includegraphics[width=0.3\textwidth]{requirement/use_case_diagram.png}
    \caption{Use Case Diagram}
\end{figure}

\subsection{Use Case Descriptions}

\begin{enumerate}[leftmargin=80pt]

\item \textbf{Detect and Visualize Sound}
    
\textbf{Actor:} Deaf User, General User

\textbf{Description:} The system captures environmental audio and displays it visually on the interface.


\item \textbf{Identify Sound Type}

\textbf{Actor:} Deaf User

\textbf{Description:} The system classifies detected sounds into predefined categories such as speech, car, or dog, and displays the result using icons or labels.

\item \textbf{View Estimated Sound Direction}

\textbf{Actor:} Deaf User

\textbf{Description:} The system estimates the direction of a sound source and represents it visually using spatial positioning on the interface.

\item \textbf{View Estimated Sound Distance}

\textbf{Actor:} Deaf User

\textbf{Description:} The system estimates the distance of a sound source and represents it visually using spatial positioning on the interface.

\item \textbf{Receive Real-Time Updates}

\textbf{Actor:} Deaf User

\textbf{Description:} The system continuously updates sound information to reflect changes in the environment in real time.

\item \textbf{Explore Sound Map}

\textbf{Actor:} General User, Researcher and Educator

\textbf{Description:} The user observes and interacts with the visual map to understand the distribution and characteristics of surrounding sounds.

\item \textbf{Analyze Sound Data}

\textbf{Actor:} Researcher and Educator

\textbf{Description:} The system provides tools for analyzing collected sound data to study sound behavior and interactions.

\item \textbf{Access Categorized Data}

\textbf{Actor:} Researcher and Educator

\textbf{Description:} The system provides categorized and labeled sound data for research, evaluation, or experimentation purposes.

\item \textbf{Evaluate System Usability}

\textbf{Actor:} Assistive Technology Stakeholders

\textbf{Description:} Accessibility experts evaluate the usability of the system to assess its effectiveness for target users.

\item \textbf{Check Accessibility Compliance}

\textbf{Actor:} Assistive Technology Stakeholders

\textbf{Description:} The system is assessed against accessibility standards to ensure it is usable by a wide audience.

\item \textbf{Maintain System}

\textbf{Actor:} Developer

\textbf{Description:} Developers maintain the system to ensure stability, reliability, and long-term functionality.

\item \textbf{Improve Machine Learning Model}

\textbf{Actor:} Developer

\textbf{Description:} Developers improve and update machine learning models to enhance sound classification accuracy.

\item \textbf{Optimize Performance}

\textbf{Actor:} Developer

\textbf{Description:} Developers optimize system performance to ensure efficient real-time processing, especially on mobile devices.

\item \textbf{Ensure Modularity}

\textbf{Actor:} Developer

\textbf{Description:} Developers design and maintain a modular system architecture to support scalability and ease of maintenance.

\end{enumerate}

\section{User Interface Design}
\label{section:user-interface-design}

\begin{figure}[h]
    \centering
    \includegraphics[width=0.8\textwidth]{requirement/interface_design.png}
    \caption{User Interface Design}
\end{figure}